\chapter{Libreria Estandar C++ (STD)}

\section{Queue (cola)}
Las colas son estructuras first-in-first-out (FIFO). \\
\textbf{Complejidad de todas las operaciones:} O(1)

\begin{lstlisting}[style=cppstyle, caption={Funcionamiento de la queue}]
	// Crear una queue:
	queue<int> cola;
	
	// Insertar elemento al FINAL de la cola: 
	cola.push(1)
	
	// Ver cuantos elementos tiene: 
	cola.size()
	
	// Eliminar el primer elemento: 
	cola.pop();
	
	// VER cual es el primer elemento: 
	cola.front()
	
	// Ver si la cola esta vacia: 
	if (cola.empty())
\end{lstlisting}

\section{Stack (pila)}
Las colas son estructuras last-in-first-out (LIFO). \\
\textbf{Complejidad de todas las operaciones:} O(1)

\begin{lstlisting}[style=cppstyle, caption={Funcionamiento del stack}]
	// Crear un stack:
	stack<int> pila;
	
	// Insertar elemento al PRINCIPIO del stack: 
	pila.push(1);
	
	// Ver cuantos elementos tiene: 
	pila.size()
	
	// Eliminar el primer elemento: 
	pila.pop();
	
	// VER cual es el primer elemento: 
	pila.top()
	
	// Ver si la cola esta vacia: 
	if (pila.empty())
\end{lstlisting}

\section{Set (conjunto)}
Almacena valores unicos de un determinado tipo, ordenados automaticamente.


\begin{lstlisting}[style=cppstyle, caption={Funcionamiento del stack}]
	// Crear un set:
	set<int> s;
	
	// Insertar un elemento:
	s.insert(3);  // O(logn)
	
	// Ver cuantos elementos tiene: // O(1)
	s.size()
	
	// Eliminar un elemento determinado: 
	// Devuelve cuantos elementos borro
	s.erase(1);  //  O(logn)
	
	// Borra todos los elementos
	us.clear();
	
	// Verificar si un valor esta de forma rapido: 
	if (s.count(x)) { ... }
	
	/* 
	Busca si un valor existe en el set.
	Si no lo encuentra, devuelve s.end()
	*/
	
	s.find(3)  // O(logn)
	
	auto it = s.find(10);
	if (it != s.end()) {
		cout << "Encontrado: " << *it << endl;
	} else {
		cout << "No esta el 10\n";
	}
	
	// Otra especie de busqueda
	auto it = s.lower_bound(25); // primer elemento >= 25
	auto it2 = s.upper_bound(30); // primer elemento > 30
	
	if (it != s.end()) cout << *it << endl;
	if (it2 != s.end()) cout << *it2 << endl;
	
\end{lstlisting}

\textbf{Otras alternativas a \texttt{set}:}
\begin{itemize}
	\item \texttt{multiset}: Permite duplicados, ordenados tambien.
	\item \texttt{unordered\_set}: Operaciones en $\mathcal{O}(1)$ promedio, pero sin orden.
\end{itemize}

\section{Representación Binaria en C++}
\subsection{Mostrar un número en binario}

La forma más directa es usando \texttt{bitset} de la STL:

\begin{lstlisting}[language=C++, caption=Muestra un número en binario con 15 bits]
	int main() {
		int x = 13;
		cout << bitset<15>(x) << endl; // Muestra: 000000000001101
	}
\end{lstlisting}

\subsection{Obtener la representación binaria como cadena}

Podemos obtener la representación binaria como un \texttt{string}:

\begin{lstlisting}[language=C++, caption=Convertir un número a cadena binaria]
	int main() {
		int x = 13;
		string bin = bitset<15>(x).to_string(); // "000000000001101"
	}
\end{lstlisting}

\subsection{Conversión binaria manual}

También podemos implementar la conversión binaria manualmente:

\begin{lstlisting}[language=C++, caption=Conversión manual de decimal a binario]
	
	void mostrar_binario(int n) {
		if (n == 0) {
			cout << "0";
			return;
		}
		string bin = "";
		while (n > 0) {
			bin = char('0' + (n % 2)) + bin;
			n /= 2;
		}
		cout << bin << endl;
	}
\end{lstlisting}



 